\subsection{Compactness and variation of finite-perimeter sets}

For an integrable function $u$, write $u\in BV(M)$ when its
distributional gradient $Du$ is a finite vector-valued measure, and
write $|Du|$ for its total variation. For $u=\chi_E$, this agrees
with the perimeter measure already defined.

Existence of confined minimizers will follow from compactness, and smooth
flows will change their volume while controlling full ambient perimeter.
We recall these facts in the precise forms used below. In smooth coordinate charts,
De Giorgi's reduced-boundary representation and Gauss--Green formula
\cite[Synopsis, pp.~119--120]{Maggi2012} give the Riemannian identities
\[
 D\chi_E=-N_E\,\mathcal H_g^{n-1}\lfloor\partial^*E,
 \qquad
 \int_E\divg X\,d\vol_g
 =\int_{\partial^*E}\langle X,N_E\rangle\,d\mathcal H_g^{n-1}.
\]
Here $\partial^*E$ is countably rectifiable, $N_E$ is its outward
measure-theoretic unit normal, and its tangent plane is $N_E^\perp$
almost everywhere. The formulas apply to every smooth compactly
supported field $X$.
To see the metric conversion, write $g$ for the metric matrix,
$\rho=\sqrt{\det g}$ and $\nu$ for the Euclidean outward normal.
If $a=(\nu^Tg^{-1}\nu)^{1/2}$, then
\[
 N_E=g^{-1}\nu/a,\qquad
 d\mathcal H_g^{n-1}=\rho a\,d\mathcal H_e^{n-1},\qquad
 \divg_gX=\rho^{-1}\divg(\rho X).
\]
The first two identities follow from orthogonality and the tangent
Gram determinant; substitution in Euclidean Gauss--Green gives the
displayed Riemannian formula. Smooth positivity of the metric and
density compares the BV norms on compact subcharts. Chart partitions
then give the invariant measure identity.

\begin{lemma}[Compactness and lower semicontinuity]\label{lem:bv-compactness}
Let $K$ be a compact subset of a smooth Riemannian manifold without
boundary. If functions $u_j$ vanish outside $K$ almost everywhere and
\[
 \sup_j\bigl(\|u_j\|_{L^1}+|Du_j|(M)\bigr)<\infty,
\]
then a subsequence converges in $L^1(M)$. Total variation is lower
semicontinuous under this convergence. If the $u_j$ are characteristic
functions, their limit is a characteristic function.
\end{lemma}

\begin{proof}
Choose a finite smooth partition on a neighborhood of $K$, with each
partition function compactly supported in a coordinate chart.
The smooth multiplier formula follows by inserting a multiplier in
the distributional test field. Smooth positivity of the metric and
volume density on the compact chart supports therefore bounds the
Euclidean BV norms of the localized functions. Extend each of them
by zero within its chart.

For a Euclidean BV function $v$, convolution satisfies
$D(v*\rho_\varepsilon)=Dv*\rho_\varepsilon$ and
$\|D(v*\rho_\varepsilon)\|_1\le |Dv|(\R^n)$.
Apply the fundamental theorem of calculus to the convolution and let
its scale tend to zero in $L^1$. This gives
\[
 \|v(\cdot+h)-v\|_1\le |h|\,|Dv|(\R^n),\qquad
 \|v*\rho_\varepsilon-v\|_1
 \le\varepsilon C_\rho |Dv|(\R^n).
\]
For fixed $\varepsilon$, the convolutions have common compact support
and uniformly bounded sup norms and first derivatives, by the uniform
$L^1$ bound on $v$. Approximation on a finite mesh makes this family
precompact in $L^1$. Taking scales to zero and a diagonal subsequence
gives compactness in each chart, hence for the original finite sum.
Total variation is a supremum of continuous divergence integrals, so
is lower semicontinuous. A further almost-everywhere convergent
subsequence of characteristic functions has values in $\{0,1\}$.
\end{proof}

\begin{lemma}[Perimeter under smooth flows]\label{lem:bv-flow}
Let $F:M\to M$ be an orientation-preserving smooth diffeomorphism,
and let $E$ have finite perimeter. If the right side below is finite,
then
\begin{equation}\label{eq:bv-flow-jacobian}
 P(FE)=\int_{\partial^*E}
 J_{n-1}\bigl(dF_x|_{N_E(x)^\perp}\bigr)\,d\mathcal H_g^{n-1}(x).
\end{equation}
All Jacobians and adjoints use the source and target Riemannian
metrics. For a smooth flow $F_t$ with velocity $Z$ and a set whose
perimeter support is compact,
\begin{equation}\label{eq:bv-flow-first-variation}
 \left.\frac d{dt}\right|_{t=0}P(F_tE)
 =\int_{\partial^*E}\divg_{N_E^\perp}Z\,d\mathcal H_g^{n-1}.
\end{equation}
For a fixed finite family of flows and perimeter supports contained
in a fixed compact set, there are common $C,\tau>0$ such that
$|P(F_tE)-P(E)|\le C|t|P(E)$ for $|t|<\tau$.
\end{lemma}

\begin{proof}
Write $J_F$ for the positive ambient volume Jacobian. The Piola
pullback of a target test field $X$ is
\[
 Q_X(x)=J_F(x)(dF_x)^{-1}X(Fx).
\]
Ordinary smooth change of variables, tested against smooth scalar
functions, gives
$\divg Q_X=J_F(\divg X)\circ F$. Gauss--Green for $E$ implies
\[
 \int_{FE}\divg X\,d\vol_g
 =\int_{\partial^*E}
 \langle X(Fx),J_F(x)(dF_x)^{-*}N_E(x)\rangle_g
 \,d\mathcal H_g^{n-1}(x).
\]
Thus the outward vector measure $-D\chi_{FE}$ is the pushforward of
the vector density $a(x)=J_F(x)(dF_x)^{-*}N_E(x)$ against the
perimeter measure of $E$. Since $F$ is injective and $a\ne0$, its
polar direction at $y=F(x)$ is $a(x)/|a(x)|$. Its total variation is
therefore the pushforward of $|a|\,d|D\chi_E|$, with no cancellation.
Linear algebra in orthonormal bases gives
$|a|=J_{n-1}(dF|_{N_E^\perp})$, proving
\eqref{eq:bv-flow-jacobian}. Differentiation of the Gram determinant
at $t=0$ gives tangential divergence. The Jacobians and their time
derivatives are uniformly bounded over all tangent planes above the
fixed compact set. This justifies differentiation under the integral
and proves the final estimate.
\end{proof}

The compactness statement supplies minimizers and compact families of
competitors. The flow formula counts contact perimeter as well as free
perimeter, both for volume adjustment and for one-sided variations.
